\section{HyperLogLog Parity: Reproduction Details}\label{app:classical-parity}

Section~\ref{sec:hll-parity} summarizes the HyperLogLog parity
experiment. This appendix covers the pieces a reader would need in
order to re-run the experiment, understand the learned-$g$
register-space analysis, and audit the full per-cell grid. Code, raw
outputs, and the benchmark driver live under
\texttt{scripts/run\_classical\_parity\_benchmark.py} in the companion
repository; full logged outputs are in \texttt{outputs/classical\_parity/hll/}.

\subsection{Protocol}

Every cell in the sweep is a training call with a zero-optimization
trainer that wraps a classical HLL into the tree reduction. Setting
$n_{\mathrm{leaves}} = 1$ recovers the flat reference;
$n_{\mathrm{leaves}} = L$ runs the tree reduction with the same
other knobs held constant. Two classical HLL implementations share
the protocol. The \textbf{native} path is a plain register-based
implementation whose merge is pointwise max on a \texttt{uint8}
register array, so register state is byte-deterministic across any
A1--A3-satisfying pipeline. The \textbf{DataSketches} path wraps
Apache DataSketches' HLL sketch and union: its internal representation
transitions between list, sparse, and dense modes based on cardinality,
so the serialized bytes of an equivalent summary can depend on the
build path even though the estimate does not.

\paragraph{Oracle function as a swappable knob.}
The oracle $\fstar$ is a configurable argument on the parity preset
with two standard choices. The \emph{analytic} setting uses true
distinct cardinality, $\fstar(x) = |\mathrm{set}(x)|$. The
\emph{HLL-reference} setting replaces it with the classical HLL's own
scoring head, $\fstar(x) = \alpha_m m^2 / \sum_j 2^{-r_j}$ applied to
the flat token span. The same $\fstar$ defines per-leaf (C1),
per-merge (C3), and root targets, so a learned merge $g$ trained under
HLL-reference receives exactly the classical sketch's scoring as
supervision at every node.

\subsection{Sweep Grid}

The default grid (used for Figure~\ref{fig:classical-parity-hll} and
Table~\ref{tab:classical-parity-hll}) is:

\begin{itemize}
\item \textbf{Classical methods:} backends $\{$native, DataSketches$\}$ $\times$ precisions $\{7, 9, 11, 13\}$ $\times$ leaf counts $\{1, 2, 4, 8, 16\}$ $\times$ oracle kinds $\{$analytic, HLL-reference$\}$ $\times$ $3$ seeds $= 240$ cells.
\item \textbf{Learned methods:} $\{g$-only, $f{+}g\}$ $\times$ precisions $\{7, 9, 11, 13\}$ $\times$ leaf counts $\{1, 2, 4, 8, 16\}$ $\times$ oracle kind HLL-reference $\times$ $3$ seeds $= 120$ cells.
\end{itemize}

Total: $360$ cells, run in parallel over CPU cores. On a 128-core host
the full sweep completes in a few minutes.

\subsection{Learned Variants in Detail}\label{app:classical-parity-variants}

The learned methods share the same per-leaf (C1), per-merge (C3), and root
targets supplied by the chosen oracle (HLL-reference or analytic). The
broad classical-sketch suite recognises four canonical variants, built from
two single-component primitives composed in stage order:

The architecture is uniform across variants: a learned readout $f$, a
learned merge $g$, and a small merge adapter all exist at the same
sizes. "Held fixed" means the component's weights are not updated this
stage; there is no identity-merge passthrough or fixed-sigmoid readout
edge case. A held-fixed component starts from a deterministic seeded
init (or from a caller-supplied checkpoint via
\texttt{init\_f\_from}/\texttt{init\_g\_from}) and stays there.

We use \emph{joint} as the umbrella codename for any schedule that
trains both components (i.e.\ any variant string with both \texttt{f}
and \texttt{g}); $f$ and $g$ stand alone for the single-component
schedules. So the experimental rows fall under three method
codenames: \texttt{learned\_f}, \texttt{learned\_g}, and
\texttt{learned\_joint}, with the specific schedule (e.g.\ \texttt{fg}
vs \texttt{gf} vs \texttt{fgf}) recorded as a per-row variant tag.

\begin{itemize}
\item \textbf{$f$:} one f stage. Train $f$; $g$ stays at its base.
\item \textbf{$g$:} one g stage. Train $g$; $f$ stays at its base. (In
  the HLL parity experiment specifically, the broad-suite model is
  replaced by a register-space-constrained sibling whose fixed readout
  is the classical HLL estimator formula
  $\alpha_m m^2 / \sum_j 2^{-r_j}$; the broad scalar-sketch model has
  no structured fixed readout.)
\item \textbf{joint} (canonical schedule \texttt{fg}): f stage, then g
  stage. After the f stage, freeze the trained $f$ and train $g$
  against it.
\item \textbf{joint, $g$-first} (schedule \texttt{gf}): g stage, then
  f stage. Symmetric ordering of the same joint training; useful as an
  ablation of which component leads.
\end{itemize}

The variant string is the schedule. Each letter is one
\texttt{fit()}-dispatched training stage that updates that component;
the next stage's frozen counterpart is loaded from the prior stage's
checkpoint, and a same-letter stage later in the string warm-starts
from that component's most recent prior checkpoint. So the underlying
$f$ and $g$ models accumulate improvements as the schedule runs.
Longer schedules such as \texttt{fgf} or \texttt{fgfgf} (a multi-pass
joint) compose the same way.

All four variants train with the oracle-plus-projection loss of
Equation~\eqref{eq:oracle-projection-objective-generalized} at
$\lambda = 0.3$ and uniform $\rho_i = 1$, with $150$ epochs per stage
over $128$ training documents, CPU-only (so the staged variants take
twice the wall time of the single-stage ones for the same per-stage
budget). Either $f$ or $g$ may additionally be \emph{supplied} as a
pretrained submodule via \texttt{init\_f\_from} / \texttt{init\_g\_from};
when supplied, it warm-starts the corresponding stage instead of
defaulting to fresh random init.

The variant string is interpreted character-by-character, so longer
sequences such as $fgf$ (``train $f$, then $g$, then $f$ again'') compose
naturally and share the same staged-checkpoint chaining. Joint
single-call training (one optimizer over both components) is not used.

\subsection{The Register-Space Constraint in Detail}

The headline scaling divergence (Section~\ref{sec:hll-parity}) deserves
more quantitative detail. At precision $p = 13$:

\begin{itemize}
\item \texttt{learned\_gf} improves with $L$: relative MAE $\approx 0.27$ at $L=2$, $\approx 0.19$ at $L=8$, $\approx 0.19$ at $L=16$. Deeper trees give more gradient signal and the learned $f$ can compensate for imprecision anywhere in the state.
\item \texttt{learned\_g} degrades with $L$: relative MAE $\approx 0.27$ at $L=1$, $\approx 0.43$ at $L=2$, $\approx 0.78$ at $L=8$, $\approx 0.86$ at $L=16$. Each merge's error channels through a fixed non-learnable readout, so mismatches at internal nodes compound rather than average out.
\end{itemize}

The two variants are trained on identical data with identical
hyperparameters; the only structural difference is whether the readout
is learnable. The gap and its direction-of-scaling with $L$ is therefore
a clean statement: some learned-$g$ settings pay a real cost when the
readout must live in the classical sketch's exact state space.

\subsection{Training-Budget Verification}

We ran the sweep at both $30$ and $150$ epochs to verify that the
\texttt{learned\_g} plateau is structural rather than a budget artifact.
At $p=13$, \texttt{learned\_g}'s relative MAE at $L=16$ moves from
$\approx 0.86$ (30 epochs) to $\approx 0.89$ (150 epochs); at $L=8$ from
$\approx 0.78$ to $\approx 0.72$; at $L=2$ from $\approx 0.43$ to
$\approx 0.50$. The magnitudes bounce within single-seed noise across
the two budgets; the direction of scaling with $L$ is identical.
\texttt{learned\_gf} shows the mirror-image stability: at $p=13, L=8$ it
holds $\approx 0.18$--$0.19$ across both budgets.

The register-space constraint is therefore not a training-rate issue:
$5\times$ the training budget does not close the gap between
\texttt{learned\_g} and \texttt{learned\_gf}. The structural point ---
that a $g$ whose outputs must live in the classical sketch's exact state
space pays a real cost that compounds with tree depth --- is a
converged observation, not a convergence artifact. A practitioner
building on this framework who does not need byte-for-byte parity with
a specific reference sketch should pick the end-to-end learned path.

\subsection{Extension to Other Classical Sketches}

A broader comparison suite beyond the HLL reproduction table compares
native HLL and Apache DataSketches HLL, CPC, Theta, Count-Min, Frequent
Items, KLL, classic quantiles, REQ, and t-digest against exact controls
and a deliberately wrong sum-of-leaf-uniques baseline. Tuple and VarOpt
supply the sampling/tuple family. The default grid mirrors the HLL
parity grid structure: capacities
\(\{\text{small},\text{medium},\text{large}\}\) crossed with leaf counts
\(L\in\{1,2,4,8,16\}\) and seeds. The report emits memory/error curves,
merge-schedule spread, empirical bound coverage, and distance to the
best official DataSketches row for each query family.
This broad grid is also routed through the unified
\texttt{fit()} abstraction as a zero-optimization trainer, so the HLL
parity rows, the multi-sketch official baselines, and learned companions
share one orchestration path. For scalar sketch readouts where an
official state is not differentiable or not the desired representation,
the companion learned-\(f{+}g\) task trains both the latent merge state
and the scalar readout jointly against either exact targets or official
sketch query functions. The same broad grid also includes a learned-\(g\)
variant with a fixed non-trainable scalar readout, isolating how much of
the performance comes from learning a locally compositional merge state
rather than from adapting the readout.

We therefore separate three statuses. \emph{Theorem-backed} rows have a
formal local-law artifact (tracked in Appendix~\ref{app:proof-artifacts})
and a runtime adapter. \emph{Official empirical} rows wrap the Apache
DataSketches implementation and are checked by the benchmark, but are
not presented as formal theorems for that exact library implementation.
\emph{Control} rows are exact references or negative baselines. Current
theorem-backed local-law bundles cover HLL, Count-Min, KLL quantiles,
GK's sequential quantile surface, and the paper's bigram/Markov-count
sketches. CPC, Theta/KMV, Frequent Items, classic quantiles, REQ,
t-digest, Tuple, and VarOpt are empirical-only until corresponding
formal artifacts are added.

\begin{table}[h]
    \centering
    \small
    \begin{tabular}{p{0.28\linewidth}p{0.22\linewidth}p{0.40\linewidth}}
        \toprule
        Sketch family & Status in this paper & Runtime comparison \\
        \midrule
        Native HLL & theorem-backed & In-repo register sketch; byte-exact tree/flat check \\
        DataSketches HLL & official empirical & Apache DataSketches HLL estimate-equivalence check \\
        Count-Min, KLL & theorem-backed + official empirical & Formal local-law artifacts plus DataSketches adapters \\
        GK & theorem-backed & Sequential-merge theorem only; no DataSketches runtime row \\
        CPC, Theta/KMV, Frequent Items, classic quantiles, REQ, t-digest & official empirical & DataSketches adapters; no formal theorem claimed \\
        Tuple, VarOpt & official empirical & DataSketches tuple/sampling adapters; no formal theorem claimed \\
        exact set & control & Exact ground truth for distinct counting \\
        sum of leaf uniques & negative control & Deliberately wrong non-mergeable distinct-count baseline \\
        \bottomrule
    \end{tabular}
    \caption{Implementation status used by the broad classical-sketch comparison suite.}
    \label{tab:classical-sketch-status}
\end{table}

\begin{figure}[t]
    \centering
    \includegraphics[width=\linewidth]{classical_sketches_summary.png}
    \caption{Broad classical-sketch comparison suite by sketch family. Each
    curve reports the best large-capacity row for one method class, averaged
    across the query heads in the family. The official floor is flat at zero
    for exact or near-exact families; the learned curves show how the learned
    merge state improves as the tree supplies more local-law supervision.}
    \label{fig:classical-sketches-summary}
\end{figure}

\begin{figure}[t]
    \centering
    \includegraphics[width=\linewidth]{classical_sketches_gold_gap.png}
    \caption{Learned rows shown as excess RMSE above the official-sketch floor.
    The vertical axis is \(\mathrm{RMSE}_{\mathrm{learned}} -
    \min_{\mathrm{official}}\mathrm{RMSE}_{\mathrm{official}}\), so zero is
    the best official DataSketches or in-repo sketch available for that query.
    Distinct counting and set union show decreasing excess error as \(L\)
    increases, while quantile heads are already close to the official floor.}
    \label{fig:classical-sketches-gold-gap}
\end{figure}

\begin{figure}[p]
    \centering
    \includegraphics[width=\linewidth]{classical_sketches_method_official.png}
    \caption{Raw, unnormalized error curves for the official empirical sketch
    rows. Each panel is a query head and each curve is a capacity preset.}
    \label{fig:classical-sketches-method-official}
\end{figure}

\begin{figure}[p]
    \centering
    \includegraphics[width=\linewidth]{classical_sketches_method_learned_f.png}
    \caption{Raw, unnormalized error curves for the learned-\(f\) rows with a
    deterministic identity merge. This isolates the contribution of the
    learned readout when the merge state is not learned.}
    \label{fig:classical-sketches-method-learned-f}
\end{figure}

\begin{figure}[p]
    \centering
    \includegraphics[width=\linewidth]{classical_sketches_method_learned_g.png}
    \caption{Raw, unnormalized error curves for the learned-\(g\) rows with a
    fixed scalar readout. This isolates the contribution of the learned
    locally compositional merge state.}
    \label{fig:classical-sketches-method-learned-g}
\end{figure}

\begin{figure}[p]
    \centering
    \includegraphics[width=\linewidth]{classical_sketches_method_learned_joint.png}
    \caption{Raw, unnormalized error curves for the learned-joint rows.
    Multi-letter schedules (\texttt{fg}, \texttt{gf}, longer) all roll up
    to the joint codename; the per-row \texttt{learned\_variant} field
    in the aggregate CSV records the specific schedule.}
    \label{fig:classical-sketches-method-learned-joint}
\end{figure}

\begin{table}[t]
    \centering
    \small
    % Auto-generated by treepo.bench.reports.classical_sketches; do not edit.
\begin{tabular}{lllrrrrrrrr}
\toprule
family & query & best official & official & learned $f$ & $f$ RMSE & learned $g$ & $g$ RMSE & learned joint (best) & joint variant & joint RMSE \\
\midrule
distinct & cardinality & theta\_datasketches & 0 & -- & -- & learned\_g\_exact\_distinct\_union\_state\_space & 0 & learned\_joint\_exact\_distinct\_union\_state\_space & fg & 0 \\
frequency & focus\_frequency & -- & -- & -- & -- & learned\_g\_count\_min\_state\_space & 0 & learned\_joint\_count\_min\_state\_space & fg & 0 \\
frequency & top5\_point\_frequency & count\_min\_datasketches & 0 & -- & -- & -- & -- & -- & -- & -- \\
quantile & rank\_at\_q0.5 & quantiles\_floats\_datasketches & 0.002583 & -- & -- & -- & -- & learned\_joint\_quantiles\_reference\_q0.5 & fg & 0.02422 \\
quantile & rank\_at\_q0.95 & tdigest\_double\_datasketches & 0.002462 & -- & -- & -- & -- & learned\_joint\_kll\_reference\_q0.95 & fg & 0.007558 \\
sampling & accumulator\_summary\_sum & tuple\_accumulator\_datasketches & 0 & -- & -- & -- & -- & learned\_joint\_tuple\_summary\_sum\_reference & fg & 0.007204 \\
sampling & total\_weight & varopt\_strings\_datasketches & 0 & -- & -- & learned\_g\_exact\_total\_weight\_state\_space & 0 & learned\_joint\_exact\_total\_weight\_state\_space & fg & 0 \\
set & a\_not\_b & theta\_datasketches & 0 & -- & -- & -- & -- & learned\_joint\_exact\_set\_a\_not\_b & fg & 0.09717 \\
set & intersection & theta\_datasketches & 0 & -- & -- & -- & -- & learned\_joint\_exact\_set\_intersection & fg & 0.2417 \\
set & union & theta\_datasketches & 0 & -- & -- & -- & -- & learned\_joint\_exact\_set\_union & fg & 0.03052 \\
\bottomrule
\end{tabular}

    \caption{Compact learned-overlay summary for the broad classical-sketch
    grid. For each query family, the table reports the best large-capacity
    official row and the best large-capacity learned-\(g\) and
    learned-joint companions.}
    \label{tab:classical-sketches-compact}
\end{table}

Conceptually, C-TreePO reduces to these classical sketches whenever the
task oracle \(\fstar\) is set equal to the official sketch query function
and the leaf/merge state is constrained to the official sketch state.
The HLL register-space experiment is the tightest reproduction of that
collapse; the broader suite supplies the same comparison harness across
the other standard mergeable-sketch families.
